\documentclass[11pt]{amsart}


\usepackage{geometry}                % See geometry.pdf to learn the layout options. There are lots.
\geometry{letterpaper}                   % ... or a4paper or a5paper or ... 

\usepackage{url}    % Need this package to display URLs properly.
\usepackage{xcolor}  % Use this package to use colored text.
\usepackage{epsfig} % Use this package to insert figures created externally.

\newcommand{\bi}{\bigskip}  % You can define your own commands to save yourself writing.

\title{An introduction to \LaTeX~for students}
\author{Christopher Hanusa}
\date{\today}


%%%%%%%%%%%%%%%%%%%%%%%%%%
\begin{document}
%%%%%%%%%%%%%%%%%%%%%%%%%%
\maketitle
\thispagestyle{empty} % Remove page numbers.


%%%%%%%%%%%%%%%%%%%%%%%%%%
\section{Introduction}

Here is some text.  We now discuss the background in Section~\ref{sec:background}.  Then we will discuss my favorite equations in Section~\ref{sec:eqns}.

%%%%%%%%%%%%%%%%%%%%%%%%%%
\subsection{Background}
\label{sec:background}

Here is the necessary background information.  Blah blah blah.  Blah blah blah.  Blah blah blah.

%%%%%%%%%%%%%%%%%%%%%%%%%%
\subsection{My Favorite Equations}
\label{sec:eqns}

\begin{equation}
\int_{1}^{9} \frac{\partial}{\partial y}\Big(yx^5+e^{xy}\Big)\, dx
\label{eqn:int}
\end{equation}

\begin{equation}
\sum_{n=1}^{\infty} \frac{1}{n^2} =\frac{\pi}{6}
\label{eqn:sum}
\end{equation}

\begin{equation}
\sqrt{a^{2} + b^{2}}
\label{eqn:sqrt}
\end{equation}

\begin{equation}
\alpha\beta\gamma\delta\cdots\omega\quad \mathbb{ABC\cdots Z}\quad \mathcal{ABC\cdots Z}
\label{eqn:labels}
\end{equation}

\begin{equation}
(A \cap B) \cup C = (A \cup C) \cap (B \cup C)
\label{eqn:cups}
\end{equation}

\begin{equation}
\begin{pmatrix} 1 & 2 \\ 3 & 4 \end{pmatrix}
\begin{pmatrix} -1 & 0 \\ 3 & 1 \end{pmatrix}=
\begin{pmatrix} 5 & 2 \\ 9 & 4 \end{pmatrix}
\label{eqn:mtx}
\end{equation}

I can now reference Equations~\eqref{eqn:int} through \eqref{eqn:mtx}.

%%%%%%%%%%%%%%%%%%%%%%%%%%
\end{document} 
%%%%%%%%%%%%%%%%%%%%%%%%%%
